% บทที่ 1 บทนำ -- ที่มาและความสำคัญ วัตถุประสงค์ ขอบเขต ข้อจำกัด จริยธรรม
% ย้ายมาจากเค้าโครง: ../chapter/02-rationale.tex, 03-objectives.tex, 05-scope.tex
% ตัวเลขกลุ่มตัวอย่างและช่วงวันเก็บข้อมูลอ่านจาก codingpy/data/ (ไม่ใช่ตัวเลขที่วางแผนไว้ในเค้าโครง)
\chapter{บทนำ}

\section{ที่มาและความสำคัญ}

นักศึกษาใช้โซเชียลมีเดียทุกวัน ทั้งเรียน ติดตามข่าว และคุยกับเพื่อน
เกือบทั้งหมดผ่านมือถือ

ถ้าใช้มากเกินไปในช่วงที่นอนไม่พอและมีความกดดันจากการเรียนอยู่แล้ว
สุขภาพจิตอาจแย่ลงโดยที่หลายคนไม่ทันรู้ตัว
ความเครียดที่สะสมทีละน้อย การนอนไม่พอติดต่อกัน
หรือการเห็นชีวิตของคนอื่นแล้วเผลอเอามาเทียบกับตัวเอง ค่อย~ๆ บั่นทอนสภาพจิตใจ
คนส่วนใหญ่รู้ตัวช้า เพราะสัญญาณเตือนช่วงแรกดูเหมือนเรื่องธรรมดา
เช่น นอนไม่ค่อยหลับ สมาธิสั้นลง หรืออารมณ์ขึ้นลงง่าย

ประเภทเนื้อหาที่เสพก็อาจสำคัญพอ~ๆ กับเวลาที่ใช้
การไล่อ่านข่าวหนัก~ๆ อย่างการเมือง อาชญากรรม หรือภัยพิบัติต่อเนื่องเป็นเวลานาน
มีชื่อเรียกในงานวิจัยต่างประเทศว่า \textit{doomscrolling}
และพบว่าสัมพันธ์กับความวิตกกังวลที่สูงขึ้น \cite{dominguez2025climate}

ระหว่างเตรียมโครงงาน คณะผู้จัดทำได้สำรวจชุดข้อมูลสาธารณะที่เกี่ยวข้องบนเว็บไซต์ \mbox{Kaggle}
\cite{singh_student_social_media,singh_social_media_addiction}
และพบว่าชุดข้อมูลที่ตรวจเป็นข้อมูลสังเคราะห์ที่สร้างไว้เพื่อการฝึกฝน ไม่ได้เก็บจากผู้ตอบจริง
จึงใช้อธิบายสถานการณ์ของนักศึกษาไม่ได้ คณะผู้จัดทำจึงเก็บข้อมูลเองจากนักศึกษาจริง

โครงงานนี้ศึกษาว่าเวลาที่ใช้โซเชียลมีเดียและประเภทเนื้อหาที่เสพ
สัมพันธ์กับสุขภาวะทางจิตของนักศึกษาเพียงใด
เพื่อให้รู้ว่าควรสังเกตสัญญาณเตือนจากพฤติกรรมด้านใดเป็นพิเศษ

\section{วัตถุประสงค์ของโครงงาน}

\begin{subitems}
    \item เพื่อศึกษาว่าเวลาที่ใช้กับโซเชียลมีเดียสัมพันธ์กับสุขภาวะทางจิตของนักศึกษาหรือไม่ และมากน้อยเพียงใด
    \item เพื่อศึกษาว่าประเภทเนื้อหาที่เสพ โดยเฉพาะการติดตามข่าวเชิงลบ สัมพันธ์กับสุขภาวะทางจิตแตกต่างกันหรือไม่
    \item เพื่อศึกษาว่าความสัมพันธ์ดังกล่าวยังคงอยู่หรือไม่ เมื่อควบคุมปัจจัยด้านการนอน การออกกำลังกาย ความกดดันจากการเรียน ผลการเรียน และความเพียงพอทางการเงิน
\end{subitems}

\section{ขอบเขตของการศึกษา}

ตารางที่~\ref{tab:scope} สรุปขอบเขตของโครงงานแยกตามขั้นตอนของกระบวนการวิทยาการข้อมูล

\begin{table}[H]
\centering
\caption{ขอบเขตของการศึกษาแยกตามขั้นตอน}
\label{tab:scope}
\begin{tabularx}{\textwidth}{|l|L|}
\hline
\textbf{ขั้นตอน} & \textbf{ขอบเขต} \\
\hline
การตั้งคำถาม & ตอบคำถามวิจัย 7 ข้อในบทที่ 3 เน้นความสัมพันธ์ระหว่างพฤติกรรมกับสุขภาวะทางจิต ไม่ครอบคลุมการหาสาเหตุ \\
\hline
การเก็บข้อมูล & นักศึกษาระดับปริญญาตรี มหาวิทยาลัยขอนแก่น ตอบแบบสอบถามออนไลน์ 121 คน ระหว่างวันที่ 23 สิงหาคม ถึง 3 กันยายน 2569 รวม 12 วัน และใช้วิเคราะห์ 107 คนที่ผ่านข้อดักความใส่ใจ \\
\hline
การสำรวจข้อมูล & สถิติเชิงพรรณนา การวิเคราะห์ตัวแปรเดี่ยวและหลายตัวแปร พร้อมการตรวจสอบคุณภาพและที่มาของข้อมูล \\
\hline
การวิเคราะห์เชิงสถิติ & การวิเคราะห์สหสัมพันธ์ การเปรียบเทียบค่าเฉลี่ยระหว่างกลุ่ม และการถดถอยพหุคูณโดยควบคุมปัจจัยด้านการนอน การออกกำลังกาย ความกดดันจากการเรียน ผลการเรียน และความเพียงพอทางการเงิน ไม่ครอบคลุมการสร้างแบบจำลองการเรียนรู้ของเครื่อง \\
\hline
ผลลัพธ์ & รายงานผลการวิเคราะห์และการอภิปราย ไม่ได้สร้างระบบหรือเครื่องมือ \\
\hline
\end{tabularx}
\end{table}

\section{ข้อจำกัดของการศึกษา}

\begin{subitems}
    \item ผลที่ได้เป็นความสัมพันธ์ ไม่ใช่เหตุและผล เพราะเก็บข้อมูลครั้งเดียว จึงบอกไม่ได้ว่าโซเชียลมีเดียทำให้สุขภาวะแย่ลง หรือผู้ที่สุขภาวะไม่ดีอยู่แล้วใช้โซเชียลมีเดียมากกว่า
    \item กลุ่มตัวอย่างเลือกแบบตามสะดวก ไม่ได้สุ่ม จึงอ้างอิงไปยังนักศึกษาทั้งมหาวิทยาลัยหรือทั้งประเทศไม่ได้
    \item ข้อมูลบางส่วนเป็นการรายงานตนเอง คะแนนในสเกล 1 ถึง 10 ของแต่ละคนอาจไม่ได้หมายถึงระดับความรู้สึกเท่ากัน ข้อชั่วโมงการใช้งานลดปัญหานี้ด้วยการให้อ่านค่าจากเครื่อง แต่ข้ออื่นยังมีปัญหานี้อยู่
    \item กลุ่มตัวอย่างมีขนาดจำกัด ถ้าความสัมพันธ์จริงมีขนาดเล็ก การศึกษานี้อาจตรวจไม่พบ โดยเฉพาะในสมการถดถอยที่มีตัวแปรต้น 8 ตัว
    \item ผลการเรียนและความเพียงพอทางการเงินเป็นข้อมูลที่ผู้ตอบรายงานเองและไม่บังคับตอบ ผู้ตอบ 3 จาก 107 คนไม่ได้ให้เกรดเฉลี่ยสะสม และคำตอบทั้งสองข้ออาจมีอคติจากการตอบให้ตัวเองดูดี
    \item ผลการศึกษาใช้ได้ในระดับกลุ่ม จึงใช้ประเมินหรือวินิจฉัยสุขภาพจิตรายบุคคลไม่ได้ ผู้ที่มีข้อกังวลควรปรึกษาผู้เชี่ยวชาญ หรือโทรสายด่วนสุขภาพจิต 1323
\end{subitems}

\section{ประเด็นด้านจริยธรรม}

\begin{subitems}
    \item ผู้ตอบต้องอ่านคำชี้แจงและกดยืนยันความยินยอมก่อนเข้าสู่แบบสอบถาม
    \item แบบสอบถามไม่เก็บชื่อ รหัสนักศึกษา อีเมล หรือข้อมูลใดที่ระบุตัวผู้ตอบได้
    \item ผู้ตอบหยุดตอบได้ตลอดเวลาโดยไม่มีผลกระทบใด~ๆ
    \item แบบสอบถามแสดงสายด่วนสุขภาพจิต 1323 ทั้งตอนต้นและตอนท้าย
    \item แบบวัด DASS-21 ไม่มีข้อคำถามเกี่ยวกับการทำร้ายตนเอง และแบบสอบถามระบุว่าผลเป็นเพียงการคัดกรองเบื้องต้น ไม่ใช่การวินิจฉัย
    \item รายงานข้อมูลในภาพรวมเท่านั้น และใช้เพื่อการศึกษาในรายวิชานี้
\end{subitems}
